\documentclass[11pt,a4paper]{article}
\usepackage{fontspec,amsmath,amssymb,amsthm,microtype}
\setmainfont{lmroman10-regular.otf}[BoldFont=lmroman10-bold.otf,ItalicFont=lmroman10-italic.otf,BoldItalicFont=lmroman10-bolditalic.otf]
\setmonofont{lmmono10-regular.otf}[Scale=MatchLowercase]
\usepackage[margin=25mm]{geometry}
\usepackage[unicode,hidelinks]{hyperref}
\setlength{\emergencystretch}{3em}
\setlength{\parskip}{0.4em}
\setlength{\parindent}{0pt}
\allowdisplaybreaks[2]
\raggedbottom
\title{\texorpdfstring{Correcting the Hessian condition-\allowbreak{}number refinement}{Correcting the Hessian condition-number refinement}}
\author{LangevinLean mathematical exposition}
\date{}
\begin{document}\maketitle
\noindent Source module: \texttt{Kinetic\allowbreak{}Langevin/\allowbreak{}Condition\allowbreak{}Number.\allowbreak{}lean}. The accompanying HTML pairs this mathematical exposition with the exact Lean statements and records the scope of the formalization.


The supplied manuscript informally simplifies its tail-free KLMC2 step/\allowbreak{}count budget using \(\kappa_2=(M_2^{2/3}+Mp^{-1/3})/m\). This module proves a valid refinement with \(\max(\kappa,\kappa_2)\), identifies the extra regime in which the printed shape follows, and constructs an exact natural-dimension family disproving a uniform removal of the ordinary condition number. Every scalar quantity is tied to the actual source definition. These are statements about the explicit budget prescription, with its Gaussian tail intentionally excluded; the separate actual-sampler accuracy theorems include that tail.


\section{\texorpdfstring{Keep both condition numbers and normalize the actual source definition}{Keep both condition numbers and normalize the actual source definition}}


The ordinary gradient condition number is \(\kappa=M/m\). The supplied source proposes \(\rho=(M_2^{2/3}+Mp^{-1/3})/m\). Introduce the dimensionless Hessian ratio \(\eta=M_2/(m\sqrt m)\). The first proof establishes the exact identity \(\rho=\eta^{2/3}+\kappa p^{-1/3}\), rather than assuming a convenient relationship between the two quantities. Positivity of \(m\) permits the real-power product and quotient rules: \((m\sqrt m)^{2/3}=m\). Lean rewrites the square root as a half power, combines exponents, distributes the quotient power, and closes the remaining fraction identity with \texttt{ring}.



{\raggedright\noindent Lean declarations: \texttt{kinetic\allowbreak{}Condition\allowbreak{}Number}, \texttt{source\allowbreak{}Hessian\allowbreak{}Condition\allowbreak{}Number}, \texttt{normalized\allowbreak{}Hessian\allowbreak{}Condition\allowbreak{}Number}, \texttt{source\allowbreak{}Hessian\allowbreak{}Condition\allowbreak{}Number\_\allowbreak{}normalize}.\par}

\section{\texorpdfstring{Derive the useful cubic inequalities from the definition}{Derive the useful cubic inequalities from the definition}}


For \(\kappa>0\), \(\eta\ge0\), and \(p>0\), the source quantity \(\rho\) is positive. Each nonnegative summand is at most their sum. Cubing \(\eta^{2/3}\le\rho\) gives \(\eta^2\le\rho^3\). Cubing \(\kappa p^{-1/3}\le\rho\), and using \((p^{-1/3})^3=1/p\), gives \(\kappa^3\le p\rho^3\). These estimates justify the later dimension scaling; they do not imply \(\kappa\le\rho\). The proof uses monotonicity of natural powers on nonnegative numbers, then multiplies through the positive dimension.



{\raggedright\noindent Lean declarations: \texttt{normalized\allowbreak{}Hessian\allowbreak{}Condition\allowbreak{}Number\_\allowbreak{}cubes}.\par}

\section{\texorpdfstring{Record the exact reciprocal-\allowbreak{}step coefficient before simplifying it}{Record the exact reciprocal-step coefficient before simplifying it}}


Choose \(\gamma=\sqrt{m+M}\), put \(s=h\sqrt m\), and target error \(\varepsilon/\sqrt m\). The two stability restrictions and the source quadratic-bias allocation give

\[s^{-1}=\max\left\{5\kappa\sqrt{1+\kappa},\ 4\sqrt{5p}\eta,\ \sqrt{\frac{4\eta p+2\sqrt2\kappa\sqrt\kappa\sqrt p}{\varepsilon}}\right\}.\]

The transient count has coefficient \(4\sqrt{1+\kappa}/s\), before its logarithm and integer rounding. \texttt{second\allowbreak{}Order\allowbreak{}Normalized\allowbreak{}Work} keeps this entire expression, including the constants. The separate \texttt{second\allowbreak{}Order\allowbreak{}Refinement} records the source's proposed shape \(\sqrt p\rho^2+\sqrt{p/\varepsilon}\rho^{5/4}\). This section concerns the source's expressly tail-free comparison; the actual sampler's separate epsilon theorem controls the Gaussian tail as well.



{\raggedright\noindent Lean declarations: \texttt{second\allowbreak{}Order\allowbreak{}Normalized\allowbreak{}Work}, \texttt{second\allowbreak{}Order\allowbreak{}Refinement}.\par}

\section{\texorpdfstring{Bound the three branches with one explicit universal constant}{Bound the three branches with one explicit universal constant}}


Suppose \(1\le\kappa\le R\), \(p\ge1\), and \(\eta^2\le R^3\). Comparing nonnegative squares gives \(\eta\le R\sqrt R\), while \(\sqrt p\le p\) and \(\kappa\sqrt\kappa\le R\sqrt R\). Thus the numerator of the accuracy branch is at most \(8pR\sqrt R\). Let \(A=\sqrt pR^2\) and \(B=\sqrt{p/\varepsilon}R^{5/4}\). Their exact squares are \(A^2=pR^4\) and \(B^2=(p/\varepsilon)R^2\sqrt R\).



After multiplying by the common count factor, the ordinary stability branch is at most \(40A\). The Hessian stability branch has squared value \(1280p(1+\kappa)\eta^2\le2560pR^4\le(64A)^2\). The accuracy branch has square at most \(256(p/\varepsilon)R^2\sqrt R=(16B)^2\). All quantities are nonnegative, so the square comparisons recover the original inequalities. Taking their maximum gives the checked bound \(64(A+B)\). The Lean proof expands these squares, applies order-preserving multiplication, and uses the real-power identity \((R^{5/4})^2=R^2\sqrt R\).



{\raggedright\noindent Lean declarations: \texttt{second\allowbreak{}Order\allowbreak{}Normalized\allowbreak{}Work\_\allowbreak{}le}.\par}

\section{\texorpdfstring{Apply the bound to the source parameters without adding a false comparison}{Apply the bound to the source parameters without adding a false comparison}}


Set \(R=\max\{\kappa,\rho\}\), where \(\rho\) is the actual source quantity. Then \(R\ge\kappa\), and the proved cubic estimate gives \(\eta^2\le\rho^3\le R^3\). The previous theorem therefore yields

\[\mathcal W(\kappa,\eta,p,\varepsilon)\le64\left[\sqrt p\,\max(\kappa,\rho)^2+\sqrt{p/\varepsilon}\,\max(\kappa,\rho)^{5/4}\right].\]

The physical-parameter wrapper substitutes \(\kappa=M/m\), \(\eta=M_2/(m\sqrt m)\), and the literal source definition of \(\rho\). Its assumptions are positive strong convexity, \(m\le M\), a nonnegative Hessian Lipschitz bound, positive accuracy, and \(p\ge1\). No comparison between the condition numbers is hidden in those assumptions.



{\raggedright\noindent Lean declarations: \texttt{second\allowbreak{}Order\allowbreak{}Normalized\allowbreak{}Work\_\allowbreak{}source\_\allowbreak{}bound}, \texttt{second\allowbreak{}Order\allowbreak{}Normalized\allowbreak{}Work\_\allowbreak{}physical\_\allowbreak{}bound}.\par}

\section{\texorpdfstring{Identify the additional regime in which the printed shape is valid}{Identify the additional regime in which the printed shape is valid}}


If the extra inequality \(\kappa\le\rho\) holds, the maximum in the corrected theorem equals \(\rho\), so the source's proposed shape is a valid upper bound with universal constant 64. The physical condition is exactly

\[M(1-p^{-1/3})\le M_2^{2/3}.\]

The equivalence follows by multiplying the original condition-number inequality by \(m>0\) and rearranging. This is a substantive extra regime, not a consequence of strong convexity or smoothness. The final wrapper makes that condition explicit and specializes the physical max theorem.



{\raggedright\noindent Lean declarations: \texttt{second\allowbreak{}Order\allowbreak{}Normalized\allowbreak{}Work\_\allowbreak{}source\_\allowbreak{}regime}, \texttt{source\allowbreak{}Hessian\allowbreak{}Condition\allowbreak{}Number\_\allowbreak{}ge\_\allowbreak{}iff}, \texttt{second\allowbreak{}Order\allowbreak{}Normalized\allowbreak{}Work\_\allowbreak{}physical\_\allowbreak{}regime}.\par}

\section{\texorpdfstring{Construct an exact family exposing the missing ordinary condition number}{Construct an exact family exposing the missing ordinary condition number}}


Take \(t>0\), \(\kappa=t^2\), \(\eta=1\), \(p=t^6\), and \(\varepsilon=1\). The exact fractional-power rules give \((t^6)^{-1/3}=t^{-2}\), so the source condition number is identically \(\rho=2\). Meanwhile the ordinary stability term alone implies \(\mathcal W\ge20\kappa^2=20t^4\). Its proposed simplification is at most \(8t^3\): \(\sqrt{t^6}=t^3\), \(2^{5/4}\le2^2=4\), and the two terms each cost at most \(4t^3\).



The proof therefore compares exact finite expressions, without relying on informal asymptotic notation. This is a counterexample to a scalar step/\allowbreak{}count simplification. It does not establish a lower bound on the possible performance of the stochastic algorithm.



{\raggedright\noindent Lean declarations: \texttt{normalized\allowbreak{}Hessian\allowbreak{}Condition\allowbreak{}Number\_\allowbreak{}family}, \texttt{second\allowbreak{}Order\allowbreak{}Normalized\allowbreak{}Work\_\allowbreak{}stability\_\allowbreak{}lower}, \texttt{second\allowbreak{}Order\allowbreak{}Refinement\_\allowbreak{}family\_\allowbreak{}upper}.\par}

\section{\texorpdfstring{Rule out every uniform multiplicative repair of the printed formula}{Rule out every uniform multiplicative repair of the printed formula}}


For any \(C>0\), choose a natural number \(n>C+1\). The preceding family has genuine integer dimension \(n^6\). Its proposed expression times \(C\) is at most \(8Cn^3\), which is strictly smaller than \(20n^4\), a lower bound on the actual reciprocal-step prescription. Hence no dimension-uniform constant \(C\) can turn the printed expression into a bound for this prescription over all admissible parameters.



The last theorem states the same contradiction with the literal physical definition \(\rho=(M_2^{2/3}+Mp^{-1/3})/m\), using \(m=M_2=1\), \(M=n^2\), and \(p=n^6\). The normalization identity identifies this physical expression with the already checked family. The conclusion concerns algebraic certification budgets; no stochastic complexity lower bound is asserted.



{\raggedright\noindent Lean declarations: \texttt{second\allowbreak{}Order\allowbreak{}Refinement\_\allowbreak{}no\_\allowbreak{}uniform\_\allowbreak{}constant}, \texttt{second\allowbreak{}Order\allowbreak{}Source\allowbreak{}Refinement\_\allowbreak{}no\_\allowbreak{}uniform\_\allowbreak{}constant}.\par}

\end{document}
